\exercisesection{}

除非另有说明，以下的希尔伯特空间均假定为复空间。

\begin{enumerate}
\def\labelenumi{\arabic{enumi})}
\tightlist
\item
  设 E 为希尔伯特空间，\(\mathcal{P}(\mathrm{E})\subset \mathcal{L}(\mathrm{E})\) 为 E 中正交投影组成的集合。设 S 为 C 的一个普遍可测子集。我们用 \(\mathcal{T}(\mathrm{S})\) 表示由满足 \(\mathrm{T}\subset \mathrm{S}\) 的 C 的普遍可测子集 T 组成的集合。S 上一个取值于 E 的投影的测度是指一个映射 \(\varpi\) （从 \(\mathcal{T}(\mathrm{S})\) 到 \(\mathcal{P}(\mathrm{E})\) ），它满足
\end{enumerate}

\begin{enumerate}
\def\labelenumi{(\roman{enumi})}
\item
  我们有 \(\varpi(\varnothing)=0\) 且 \(\varpi(S)=1_{E}\) 。
\item
  若 \((\mathrm{T}_{n})\) 是 \(\mathcal{T}(\mathrm{S})\) 中两两不相交的集合的序列，且 T 是诸集合 \(T_{n}\) 的并，则
\end{enumerate}

\[
\varpi (\mathrm{T}) = \sum_ {n \in \mathbf {N}} \varpi (\mathrm{T} _ {n})
\]

在赋有逐点收敛拓扑的空间 \(\mathcal{L}(\mathrm{E})\) 中成立。

对 E 上每个谱等于 S 的正规部分算子 u，以及对 \(T \in \mathcal{T}(S)\) ，我们用 \(p_{T,u}\) 表示由 T 定义的 u 的谱投影。

把 E 上谱为 S 的正规部分算子 \(u\) 对应于映射 \(\mathrm{T} \mapsto p_{\mathrm{T}, u}\) 的对应，是谱为 S 的正规部分算子之集与 S 上取值于 E 的投影的测度之集之间的一个双射。

\begin{enumerate}
\def\labelenumi{\arabic{enumi})}
\setcounter{enumi}{1}
\tightlist
\item
  设 E 为复希尔伯特空间，u 为 E 上的一个正规部分算子。我们用 \(p_{A}\) 表示 u 在 A 上的谱投影。设 a \textless{} b 为实数。我们有
\end{enumerate}

\[
\frac {1}{2} (p _ {[ a, b ]} + p _ {] a, b [}) = \lim _ {\varepsilon \rightarrow 0} \frac {1}{2 i \pi} \int_ {a} ^ {b} \big (\mathrm{R} (u, t - i \varepsilon) - \mathrm{R} (u, t + i \varepsilon) \big) d t,
\]

在赋有逐点收敛拓扑的空间 \(\mathcal{L}(\mathrm{E})\) 中成立（“斯通公式”）。

\begin{enumerate}
\def\labelenumi{\arabic{enumi})}
\setcounter{enumi}{2}
\item
  设 E 为可数型的希尔伯特空间。\(\mathcal{L}(\mathrm{E})\) 仅有的闭双边理想是 \(\{0\}\) 、\(\mathcal{L}^{\mathrm{c}}(\mathrm{E})\) 与 \(\mathcal{L}(\mathrm{E})\) 。（利用 IV 第 316 页习题 10；若 I 是这样一个理想，假设 I 中存在一个非紧的非零自伴元 \(u\) ，并证明：对 R 中每个使 \(0 \notin J\) 的开区间 J ，\(u\) 在 J 上的谱投影属于 I。）
\item
  设 E 为希尔伯特空间，u 为 E 上的一个自伴部分算子。映射 \(\sigma\colon f\mapsto f(u)\) （从 \(\mathcal{L}_{\mathrm{u}}^{\infty}(\mathrm{Sp}(u))\) 到 \(\mathcal{L}(\mathrm{E})\) ）是满足下列三个条件的唯一的对合代数态射：
\end{enumerate}

\begin{enumerate}
\def\labelenumi{(\roman{enumi})}
\item
  \(\| \sigma (f)\| \leqslant \| f\|_{\infty}\) 对所有 \(f\in \mathcal{L}_{\mathrm{u}}^{\infty}(\mathrm{Sp}(u))\) 成立
\item
  若 \((f_n)_{n\in \mathbf{N}}\) 是 \(\mathcal{L}_{\mathrm{u}}^{\infty}(\mathrm{Sp}(u))\) 中的有界序列，且 \(f_{n}(t)\to f(t)\) 对所有 \(t\in \mathrm{Sp}(u)\) 成立，则 \(\sigma (f_n)\rightarrow \sigma (f)\) 在赋有逐点收敛拓扑的 \(\mathcal{L}(\mathrm{E})\) 中成立。
\item
  若 \((f_n)_{n\in \mathbf{N}}\) 是 \(\mathcal{L}_{\mathrm{u}}^{\infty}(\mathrm{Sp}(u))\) 中的序列，满足 \(|f_n(t)|\leqslant |t|\) 且 \(f_{n}(t)\to t\) 对所有 \(t\in \mathrm{Sp}(u)\) 成立，则 \(u(x) = \lim \sigma (f_n)(x)\) 对所有 \(x\in \mathrm{dom}(u)\) 成立。
\end{enumerate}

\begin{enumerate}
\def\labelenumi{\arabic{enumi})}
\setcounter{enumi}{4}
\item
  设 u 为希尔伯特空间 E 上的一个正规部分算子。u 的剩余谱是空的。
\item
  设 E 与 F 为希尔伯特空间。设 u 为从 E 到 F 的一个稠定闭部分算子。设 \(\lambda \in C^{*}\) 属于 E 上自伴算子 \(u^{*} \circ u\) 的预解集（参见 IV 第 241 页命题 12）。
\end{enumerate}

\begin{enumerate}
\def\labelenumi{\alph{enumi})}
\item
  从 F 到 F 的部分算子 \(u \circ u^{*}\) 是自伴的。
\item
  对每个向量 \(x \in \mathrm{dom}(u^{*})\) ，向量
\end{enumerate}

\[
y = (1 _ {\mathrm{F}} + u \circ \mathrm{R} (u ^ {*} \circ u, \lambda) \circ u ^ {*}) x
\]

属于 dom \((u\circ u^{*})\) ，且公式

\[
(u \circ u ^ {*} - \lambda 1 _ {\mathrm{F}}) y = x
\]

成立。

\begin{enumerate}
\def\labelenumi{\alph{enumi})}
\setcounter{enumi}{2}
\tightlist
\item
  部分算子 \(u \circ \mathrm{R}(u^* \circ u, \lambda) \circ u^*\) 稠定，且在其定义域上是连续的。
\end{enumerate}

\(d)\) 部分算子 \(u \circ u^{*} - \lambda 1_{\mathrm{F}}\) 是满射的。

\begin{enumerate}
\def\labelenumi{\alph{enumi})}
\setcounter{enumi}{4}
\item
  我们有 \(\lambda \notin \operatorname{Sp}(u \circ u^{*})\) 。
\item
  我们有恒等式
\end{enumerate}

\[
\lambda \mathrm{R} (u \circ u ^ {*}, \lambda) - u \circ \mathrm{R} (u ^ {*} \circ u, \lambda) \circ u ^ {*} = 1 _ {\mathrm{F}}
\]

在 \(\mathcal{L}(\mathrm{F})\) 中成立（其中 \(u\circ \mathrm{R}(u^{*}\circ u,\lambda)\circ u^{*}\) 按记号滥用表示 F 的与此部分算子在其定义域上相重合的那个自同态，参见问题 \(c\) ））。

\begin{enumerate}
\def\labelenumi{\alph{enumi})}
\setcounter{enumi}{6}
\tightlist
\item
  我们有关系式
\end{enumerate}

\[
u ^ {*} \circ \mathrm{R} (u \circ u ^ {*}, \lambda) \subset \mathrm{R} (u ^ {*} \circ u, \lambda) \circ u ^ {*}.
\]

\begin{enumerate}
\def\labelenumi{\alph{enumi})}
\setcounter{enumi}{7}
\tightlist
\item
  部分算子 u 通过过渡到子空间，诱导出从 \(\operatorname{Ker}(u^{*} \circ u - \lambda1_{\mathrm{E}})\) 到 \(\operatorname{Ker}(u \circ u^{*} - \lambda1_{\mathrm{F}})\) 的一个同构。
\end{enumerate}

\begin{enumerate}
\def\labelenumi{\arabic{enumi})}
\setcounter{enumi}{6}
\tightlist
\item
  设 A 为希尔伯特空间 E 上一族稠定闭算子。设 \(\lambda \in \mathbf{C}\) 。我们称 A 是 \(\lambda\) -交换的，如果 \(\lambda\) 属于 \(u\) 的预解集（对所有 \(u \in \mathrm{A}\) ），且 E 的诸自同态 \(\mathrm{R}(u, \lambda)\) （\(u \in \mathrm{A}\) ）两两可交换。
\end{enumerate}

我们称 A 是预解意义上交换的，如果存在 \(\lambda \in \mathbf{C}\) 使 A 是 \(\lambda\) -交换的。

由单个稠定闭算子 u 组成的集是交换的，当且仅当 u 的预解集非空。

\begin{enumerate}
\def\labelenumi{\alph{enumi})}
\item
  若 A ⊂ L(E)，则 A 是预解意义上交换的，当且仅当 A 中各元素的谱之并不等于 C，且 A 是交换的。
\item
  设 A 为 E 上一族 λ-交换的稠定闭算子。则对属于每个 \(u \in A\) 的预解集的每个复数 μ，A 是 μ-交换的。
\item
  设 A 为 E 上一族预解意义上交换的稠定闭算子，且 \(u^{*} \in A\) 对所有 \(u \in A\) 成立。存在局部紧拓扑空间 X、X 上的正测度 \(\mu\) 、一个等距同构 \(v\) （从 \(\mathrm{L}^2 (\mathrm{X},\mu)\) 到 E ），以及映射 \(u\mapsto g_u\) （从 A 到 \(\mu\) -可测函数之集） ，使得 \(u = v\circ m_{g_u}\circ v^{-1}\) 对所有 \(u\in \mathrm{A}\) 成立。
\end{enumerate}

\begin{enumerate}
\def\labelenumi{\arabic{enumi})}
\setcounter{enumi}{7}
\tightlist
\item
  设 u 与 v 为希尔伯特空间 E 上的正规部分算子。下列断言等价：
\end{enumerate}

\begin{enumerate}
\def\labelenumi{\alph{enumi})}
\item
  自同态 \(b(u)\) 与 \(b(v)\) 可交换；
\item
  对 C 的每个普遍可测子集 A，u 与 v 在 A 上的谱投影可交换；
\end{enumerate}

若 \(u\) 与 \(v\) 的预解集非空，这些断言等价于

\begin{enumerate}
\def\labelenumi{\alph{enumi})}
\setcounter{enumi}{2}
\tightlist
\item
  对每个 \(\lambda \notin \operatorname{Sp}(u)\) 与 \(\lambda' \notin \operatorname{Sp}(v)\) ，预解式 \(\operatorname{R}(u, \lambda)\) 与 \(\operatorname{R}(v, \lambda')\) 可交换。
\end{enumerate}

若 \(u\) 与 \(v\) 自伴，这些断言等价于

\begin{enumerate}
\def\labelenumi{\alph{enumi})}
\setcounter{enumi}{3}
\tightlist
\item
  对每个 \(t \in R\) ，酉自同态 \(e^{itu}\) 与 \(e^{itv}\) 可交换。
\end{enumerate}

\begin{enumerate}
\def\labelenumi{\arabic{enumi})}
\setcounter{enumi}{8}
\tightlist
\item
  设 u 为希尔伯特空间 E 上的一个稠定闭部分算子。下列断言等价：
\end{enumerate}

\begin{enumerate}
\def\labelenumi{\alph{enumi})}
\item
  算子 u 是正规的；
\item
  存在 E 上预解意义上可交换的自伴算子 \(u_{1}\) 与 \(u_{2}\) ，使得 \(u = u_{1} + iu_{2}\) ；
\item
  我们有 \(\operatorname{dom}(u) = \operatorname{dom}(u^*)\) ，且 \(\| u(x)\| = \| u^{*}(x)\|\) 对所有 \(x\in \operatorname{dom}(u)\) 成立；
\item
  我们有 \(uu^{*} = u^{*}u\) 。
\end{enumerate}

若 \(u\) 的预解集非空，这些断言等价于

\begin{enumerate}
\def\labelenumi{\alph{enumi})}
\setcounter{enumi}{4}
\tightlist
\item
  对每个 \(\lambda \notin \operatorname{Sp}(u)\) ，预解式 \(\operatorname{R}(u, \lambda)\) 是正规的。
\end{enumerate}

\begin{enumerate}
\def\labelenumi{\arabic{enumi})}
\setcounter{enumi}{9}
\tightlist
\item
  设 E 为复希尔伯特空间，A 为 \(\mathcal{L}(\mathrm{E})\) 的一个含单位的自伴子代数，它在逐点收敛拓扑下是闭的。（这样的代数称为冯·诺伊曼代数。）
\end{enumerate}

\begin{enumerate}
\def\labelenumi{\alph{enumi})}
\item
  存在 \(\mathcal{L}(\mathrm{E})\) 的一个子集 X，使 A 与 X 在 \(\mathcal{L}(\mathrm{E})\) 中的交换子相重合。（利用 IV 第 326 页习题 24。）反之，对 X 的每个自伴子集，X 在 \(\mathcal{L}(\mathrm{E})\) 中的交换子是一个含单位的自伴子代数，且在逐点收敛拓扑下是闭的。
\item
  设 \(u \in \mathrm{A}\) ，\((j, |u|)\) 为 \(u\) 的极分解（I 第 140 页定义 4）。元素 \(j\) 与 \(u\) 属于 A。
\item
  设 \(u \in A\) 且 u 正规。对每个函数 \(f \in \mathcal{L}_{\mathrm{u}}^{\infty}(\mathrm{Sp}(u))\) ，自同态 \(f(u)\) 属于 A。
\item
  设 \(u \in A\) 。元素 u 是 A 中的右（相应地，左）零因子（参见 A，I，第 93 页）当且仅当 u 的像在 E 中不稠密（相应地，当且仅当 u 不是单射的）。
\item
  设 F 为 E 上部分算子组成的集，其元素是满足下列条件的稠定闭部分算子 u：对 X 中每个 x，部分算子 \(x \circ u\) 是部分算子 \(u \circ x\) 的一个约化。对 F 中每个 u 与 A 中每个 v，部分算子 \(u^{*}\) 、\(u + v\) 、\(v \circ u\) 与 \(u \circ v\) 都属于 F。
\end{enumerate}

\(f\) ) 设 \(u\in \mathrm{F}\) ，\((j,|u|)\) 为 \(u\) 的极分解（IV 第 290 页定义 7）。部分算子 \(|u|\) 属于 F，而 \(j\) 属于 A。

\(g)\) 设 \(u \in \mathrm{F}\) 且 \(u\) 正规。对每个函数 \(f \in \mathcal{L}_{\mathrm{u}}(\mathrm{Sp}(u))\) ，部分算子 \(f(u)\) 属于 F；若 \(f\) 有界，则 \(f(u)\) 属于 A。

\begin{enumerate}
\def\labelenumi{\alph{enumi})}
\setcounter{enumi}{7}
\item
  设 \(u \in A\) 单射且其像在 E 中稠密。部分算子 \(u^{-1}\) 属于 F。
\item
  对 F 中每个 u，存在 A 的元素 \(u_{1}\) 与 \(u_{2}\) ，使 \(u_{2}\) 单射、像稠密，且 \(u = u_{1}u_{2}^{-1}\) 。（先化归到 u 自伴且正的情形，再对函数 \(f_{1}: t \mapsto t\varphi(t) + \psi(t)\) 与 \(f_{2}: t \mapsto \varphi(t) + t^{-1}\psi(t)\) 应用普遍可测函数演算，其中 \(\varphi\) 与 \(\psi\) 分别是区间 {[}0, 1{[} 与 {[}1, +∞{[} 的特征函数。{]})
\item
  假设 A 中每个像稠密的单射元素都是满射的。则代数 A 满足 Ore 条件：对 A 中任意 u 与 s，若 s 不是零因子（既非右零因子也非左零因子），则存在 A 中 \(u'\) 与 \(s'\) ，使 \(s'\) 不是零因子且 \(us' = su'\) 。
\end{enumerate}

\(k\) ) 上一问题的结论并不总是成立。（考虑 \(\mathrm{A} = \mathcal{L}(\mathrm{E}).\) ）

\begin{enumerate}
\def\labelenumi{\arabic{enumi})}
\setcounter{enumi}{10}
\tightlist
\item
  \begin{enumerate}
  \def\labelenumii{\alph{enumii})}
  \tightlist
  \item
    设 \(q\) 为复希尔伯特空间 E 上的一个正部分形式。则 \(q\) 是闭的，当且仅当：对每个序列 \((x_{n})\) （由 \(\text{dom}(q)\) 中元素组成），若它收敛于 E 中 \(x\) 且满足
  \end{enumerate}
\end{enumerate}

\[
\lim _ {n, m \to + \infty} q (x _ {n} - x _ {m}, x _ {n} - x _ {m}) = 0,
\]

，我们有 \(x \in \text{dom}(q)\) 且 \(q(x_{n} - x, x_{n} - x) \to 0\) 。

\begin{enumerate}
\def\labelenumi{\alph{enumi})}
\setcounter{enumi}{1}
\tightlist
\item
  设 \(E = L^{2}(\mathbf{R}, \mu)\) ，其中 \(\mu\) 为勒贝格测度。以 \(\mathcal{D}(\mathbf{R})\) 为定义域、由 \(q(\varphi_{1}, \varphi_{2}) = \overline{\varphi_{1}(0)}\varphi_{2}(0)\) 定义的正埃尔米特形式 q 不是闭的，且不存在扩张 q 的闭的正部分形式。
\end{enumerate}

\begin{enumerate}
\def\labelenumi{\arabic{enumi})}
\setcounter{enumi}{11}
\tightlist
\item
  我们在 \(\mathbf{R}\) 上配备勒贝格测度。设 E 为 \(\mathrm{L}^2 (\mathbf{R})\) 中由那些 \(f\in \mathrm{L}^2 (\mathbf{R})\) （满足 \(f(-x) = f(x)\) ，对几乎所有 \(x\in \mathbf{R}\) ）组成的闭子空间。设 \(\mathcal{D}_{+}(\mathbf{R})\) 为 \(\mathbf{R}\) 上满足 \(f(x) = f(-x)\) 对所有 \(x\in \mathbf{R}\) 成立的测试函数组成的空间。我们把 \(\mathcal{D}_{+}(\mathbf{R})\) 与 E 的一个子空间等同起来。我们用 \(\mathcal{F}\in \mathcal{L}(\mathrm{E})\) 表示由 \(\mathrm{L}^2 (\mathbf{R})\) 的傅里叶变换通过过渡到子空间而导出的 E 的自同态。
\end{enumerate}

\begin{enumerate}
\def\labelenumi{\alph{enumi})}
\item
  对每个 \(f \in \mathcal{D}_{+}(\mathbf{R})\) ，我们有 \(\mathcal{F}(f) + \mathcal{F}^{*}(f) \in \mathcal{D}_{+}(\mathbf{R})\) 。
\item
  设 u 为 E 上以 \(\mathcal{D}_{+}(\mathbf{R})\) 为定义域、由 \(u(f) = \mathcal{F}(f) + \mathcal{F}^{*}(f)\) （对 \(f \in \mathcal{D}_{+}(\mathbf{R})\) ）定义的部分算子。算子 u 本质自伴。
\item
  我们有 \(\operatorname{dom}(u^{2})=\{0\}\) （利用 II 第 273 页习题 28。）
\end{enumerate}

\begin{enumerate}
\def\labelenumi{\arabic{enumi})}
\setcounter{enumi}{12}
\tightlist
\item
  设 E 为复希尔伯特空间，u 为 E 上的一个对称部分算子。假设存在实数 c \textgreater{} 0 ，使 \(\langle x \mid u(x) \rangle \geqslant c \|x\|^{2}\) 对所有 \(x \in \text{dom}(u)\) 成立。
\end{enumerate}

证明下列条件等价：

\begin{enumerate}
\def\labelenumi{(\roman{enumi})}
\item
  部分算子 \(u\) 本质自伴；
\item
  u 的像在 E 中稠密；
\item
  \(u^{*}\) 的核退缩为 0。
\end{enumerate}

\begin{enumerate}
\def\labelenumi{\arabic{enumi})}
\setcounter{enumi}{13}
\tightlist
\item
  设 \(n \in \{1, 2, 3\}\) 。我们用 \(\Delta\) 表示 \(\mathrm{L}^2(\mathbf{R}^n)\) 上的拉普拉斯算子，用 D 表示其定义域。它是一个正的部分算子。
\end{enumerate}

\begin{enumerate}
\def\labelenumi{\alph{enumi})}
\tightlist
\item
  存在实数 \(c \geqslant 0\) ，使对每个 \(f \in D\) ，我们有
\end{enumerate}

\[
\| f \| _ {\infty} \leqslant c (\| \Delta f \| + \| f \|).
\]

\begin{enumerate}
\def\labelenumi{\alph{enumi})}
\setcounter{enumi}{1}
\tightlist
\item
  对每个实数 \(a > 0\) ，存在实数 \(b \geqslant 0\) ，使对每个 \(f \in \mathrm{D}\) ，我们有
\end{enumerate}

\[
\| f \| _ {\infty} \leqslant a \| \Delta f \| + b \| f \|.
\]

（对 \(x \mapsto f(\alpha x)\) 应用上一问题，其中 \(\alpha \in \mathbf{R}_{+}^{*}\) 。）

\begin{enumerate}
\def\labelenumi{\alph{enumi})}
\setcounter{enumi}{2}
\item
  设 \(g \in \mathcal{L}^{\infty}(\mathbf{R}^{n}) + \mathcal{L}^{2}(\mathbf{R}^{n})\) 。若 \(g\) 是实值的，则部分算子 \(\Delta + m_{g}\) 自伴。
\item
  设 \(\alpha\) 为满足 \(0 \leqslant \alpha < n/2\) 的实数。设 N 为 \(\mathbf{R}^n\) 上的欧几里得范数，\(g\) 为 \(\mathbf{R}^n\) 上由 \(g(x) = 1/\mathrm{N}(x)^\alpha\) （\(x \neq 0\) 时）与 \(g(0) = 0\) 定义的函数。部分算子 \(\Delta + m_g\) 自伴。
\end{enumerate}

\begin{enumerate}
\def\labelenumi{\arabic{enumi})}
\setcounter{enumi}{14}
\tightlist
\item
  设 \(n \in \mathbf{N}\) ，\(\Delta\) 为 \(\mathrm{L}^2(\mathbf{R}^n)\) 上唯一的拉普拉斯算子。用 B 表示 \(\mathbf{R}^n\) 中的单位球。
\end{enumerate}

若 \(g \in \mathcal{L}_{\mathrm{u}}(\mathbf{R}^{n})\) 使 \(\Delta\) 的定义域包含于 \(m_{g}\) 的定义域，则

\[
\sup _ {y \in \mathbf {R} ^ {n}} \int_ {\mathrm{B}} | g (x + y) | ^ {2} d x
\]

是有限的。（首先验证 \(m_{g}\) 是 \(\Delta\) -有界的。）

\begin{enumerate}
\def\labelenumi{\arabic{enumi})}
\setcounter{enumi}{15}
\tightlist
\item
  \begin{enumerate}
  \def\labelenumii{\alph{enumii})}
  \tightlist
  \item
    设 I = {]}0, 1{[}。空间 H \(^{1}\) (I) 包含于 C(Ī)。
  \end{enumerate}
\end{enumerate}

\begin{enumerate}
\def\labelenumi{\alph{enumi})}
\setcounter{enumi}{1}
\item
  微分算子 \(u\colon f\mapsto -if'\) （\(\mathrm{L}^2 (\mathrm{I})\) 上以 \(\mathcal{D}(\mathrm{I})\) 为定义域）是对称且可闭的。对所有 \(f\in \mathrm{dom}(u^{*})\) ，导出广义函数 \(f'\) （由 \(f\) 导出）属于 \(\mathrm{L}^2 (\mathrm{I})\) ，且我们有 \(u^{*}(f) = if'\) 。
\item
  \(\overline{u}\) 的亏空间都是 1 维的；由此给出 \(\overline{u}\) 的所有自伴扩张的一个明显描述。
\end{enumerate}

17) 对 I ={]}0, 1{[} 上的拉普拉斯算子进行分类。

¶ 18) 设 n 为整数 ≥ 1。我们用 D 表示 L²(T\^{}n) 上以 C∞(T\^{}n) 为定义域、满足

\[
\mathrm{D} (\varphi) = - \sum_ {i = 1} ^ {n} \partial_ {i} ^ {2} \varphi
\]

（对所有 \(\varphi\in\mathcal{C}^{\infty}(\mathbf{T}^{n})\) ）的部分算子，

\begin{enumerate}
\def\labelenumi{\alph{enumi})}
\item
  部分算子 D 本质自伴。我们用 \(\Delta\) 表示它的闭包。
\item
  部分算子 \(\Delta\) 自伴且正。
\item
  部分算子 \(\Delta\) 有紧预解式。
\end{enumerate}

\(d)\) 设 \(\lambda\in\mathbf{R}\) 。我们有 \(\lambda\in\mathrm{Sp}(\Delta)\) 当且仅当存在自然数 \(k_{1},\ldots,k_{n}\) 使 \(\lambda=4\pi^{2}(k_{1}^{2}+\cdots+k_{n}^{2})\) 。

\begin{enumerate}
\def\labelenumi{\alph{enumi})}
\setcounter{enumi}{4}
\tightlist
\item
  设 \(\lambda\) 为 \(\Delta\) 的一个特征值；用 \(N(\lambda)\) 表示相应特征空间的维数。我们有
\end{enumerate}

\[
\sum_ {\lambda \leqslant \mathrm{T}} \mathrm{N} (\lambda) = c _ {n} \mathrm{T} ^ {n / 2} + \mathrm{O} (\mathrm{T} ^ {(n - 1) / 2})
\]

当 T → +∞ 时成立。

\(f)\) 若 \(n \geqslant 2\) ，则 \(\limsup_{\lambda} N(\lambda) = +\infty\) 。（在 \(n = 2\) 的情形，考虑 \(\lambda = 4\pi^2 m\) ，其中 \(m\) 是形如 \(4k + 1\) 的素数的乘积。） \(g)\) 若 \(n = 1\) 或 \(n = 2\) ，我们有

\[
\operatorname{Card} (\operatorname{Sp} (u) \cap [ 0, \mathrm{T} ]) = o (\mathrm{T}),
\]

并且若 \(n\geqslant 3\) ，我们有

\[
\liminf _ {T \to + \infty} \frac {1}{T} \operatorname{Card} (\operatorname{Sp} (u) \cap [ 0, T ]) > 0.
\]

\begin{enumerate}
\def\labelenumi{\arabic{enumi})}
\setcounter{enumi}{18}
\tightlist
\item
  设 u 为希尔伯特空间 E 上的一个自伴部分算子。
\end{enumerate}

\begin{enumerate}
\def\labelenumi{\alph{enumi})}
\tightlist
\item
  设 \(t \in R\) 且 \(x \in E\) 。若 x 关于 u 的谱测度具有紧支撑，则我们有
\end{enumerate}

\[
e ^ {i u} (x) = \sum_ {n \in \mathbf {N}} \frac {u ^ {n} (x)}{n !},
\]

其中级数在 E 中收敛。

\begin{enumerate}
\def\labelenumi{\alph{enumi})}
\setcounter{enumi}{1}
\tightlist
\item
  由使得 x 关于 u 的谱测度具有紧支撑的那些向量 \(x \in E\) 组成之集是 E 的一个稠密子空间。
\end{enumerate}

\begin{enumerate}
\def\labelenumi{\arabic{enumi})}
\setcounter{enumi}{19}
\tightlist
\item
  设 \(u\) 为希尔伯特空间 E 上的一个部分算子。我们用 \(\mathrm{dom}(u)^{\mathrm{an}}\) 表示由满足下列条件的向量 \(x \in \mathrm{E}\) 组成的空间：\(x \in \mathrm{dom}(u^n)\) 对每个整数 \(n \in \mathbf{N}\) 成立，且幂级数
\end{enumerate}

\[
\sum_ {n \in \mathbf {N}} \frac {\| u ^ {n} (x) \|}{n !} z ^ {n}\tag{23}
\]

的收敛半径 \textgreater{} 0。\(\text{dom}(u)^{\text{an}}\) 的元素称为 u 的解析向量。

设 \(x\) 为 E 中的一个元素，\(x\) 属于 \(\mathrm{dom}(u^n)\) （对所有 \(n \in \mathbf{N}\) ）。设 \(\mathrm{E}_x\) 为由诸向量 \(u^n(x)\) （\(n \in \mathbf{N}\) ）生成的 E 的子空间的闭包，\(u_x\) 为 \(\overline{\mathrm{E}}_x\) 上以 \(\mathrm{E}_x\) 为定义域、由 \(u\) 通过过渡到子空间导出的部分算子。我们称 \(x\) 是 \(u\) 的一个唯一性向量，如果 \(u_x\) 本质自伴。

\begin{enumerate}
\def\labelenumi{\alph{enumi})}
\tightlist
\item
  假设 u 对称，且 E 含有 u 的唯一性向量的一个完全子集。则 u 本质自伴。（证明 \(u + i1_{E}\) 的像与 \(u - i1_{E}\) 的像都在 E 中稠密。）
\end{enumerate}

以下设 u 对称，且 E 含有解析向量的一个完全子集。

\begin{enumerate}
\def\labelenumi{\alph{enumi})}
\setcounter{enumi}{1}
\item
  对 u 的每个解析向量 x，部分算子 \(u_{x}\) （在 \(\overline{E}_{x}\) 上）有一个自伴扩张。
\item
  设 \(x\) 为 \(u\) 的一个解析向量，\(r > 0\) 为幂级数 (23) 的收敛半径。设 \(v\) 为 \(u_x\) 的一个自伴扩张。对每个 \(t \in \mathbf{R}\) （满足 \(|t| < r\) ），我们有
\end{enumerate}

\[
\langle e ^ {i t v} (x) \mid x \rangle = \sum_ {n \in \mathbf {N}} \frac {(i t) ^ {n}}{n !} \langle u ^ {n} (x) \mid x \rangle .
\]

（利用 \(x\) 关于 \(v\) 的谱测度，证明左端是定义在由 \(t \in \mathbf{C}\) （满足 \(|t| < r\) ）组成之集上的一个解析函数的限制，并计算该函数在 0 处的幂级数展开。）

\begin{enumerate}
\def\labelenumi{\alph{enumi})}
\setcounter{enumi}{3}
\item
  对每个 \(t \in R\) ，酉算子 \(e^{itv}\) 只依赖于 \(u_x\) ，且算子 \(u_x\) 本质自伴。（利用 V 第 428 页定理。） e) 由此推出 u 本质自伴。
\item
  E 上闭的对称部分算子 u 自伴，当且仅当 \(\text{dom}(u)^{\text{an}}\) 在 E 中稠密。
\end{enumerate}

\begin{enumerate}
\def\labelenumi{\arabic{enumi})}
\setcounter{enumi}{20}
\tightlist
\item
  我们用 \(x\) 表示 \(\mathbf{R}\) 的恒等函数。我们用 \(\mathcal{M}_{\mathrm{m}}(\mathbf{R})\) 表示由正测度 \(\nu\) （在 \(\mathbf{R}\) 上）组成的、满足 \(x^n \in \mathcal{L}^1(\mathbf{R}, \nu)\) （对每个整数 \(n \in \mathbf{N}\) ）的集。 \(a)\) 设 \(\nu \in \mathcal{M}_{\mathrm{m}}(\mathbf{R})\) ，并令 \(\mu_n = \nu(x^n)\) （对所有 \(n \in \mathbf{N}\) ）（“测度 \(\nu\) 的矩”）。对每个整数 \(n \geqslant 0\) 与所有复数 \(a_0, \ldots, a_n\) ，我们有
\end{enumerate}

\[
\sum_ {i = 0} ^ {n} \sum_ {j = 0} ^ {n} \overline {{a}} _ {i} a _ {j} \mu_ {i + j} \geqslant 0.\tag{24}
\]

在下文中，我们打算研究问题 \(a\) ) 的逆命题（“汉堡矩问题”）。因此，在本习题的其余部分，我们给定一个满足条件 (24) 的实数序列 \(\boldsymbol{\mu} = (\mu_n)_{n \geqslant 0}\) 。

\begin{enumerate}
\def\labelenumi{\alph{enumi})}
\setcounter{enumi}{1}
\tightlist
\item
  存在唯一的正半双线性形式 \(b_{\mu}\) （在 C{[}X{]} 上），使对所有整数 i 与 j 有 \(b_{\mu}(X^{i}, X^{j}) = \mu_{i+j}\) 。
\end{enumerate}

我们用 \(E_{\mu}\) 表示 C{[}X{]} 关于形式 \(b_{\mu}\) 的分离完备希尔伯特空间。 c) 线性映射 \(u_{\mu}\) （从 C{[}X{]} 到 C{[}X{]} ，由 \(u_{\mu}(P) = XP\) 定义）通过过渡到商，诱导出一个对称部分算子 \(u_{\mu}\) ，位于希尔伯特空间 \(E_{\mu}\) 上。

\begin{enumerate}
\def\labelenumi{\alph{enumi})}
\setcounter{enumi}{3}
\item
  部分算子 \(u_{\mu}\) 有一个自伴扩张。
\item
  存在测度 \(\nu\in\mathcal{M}_{\mathrm{m}}(\mathbf{R})\) 使 \(\mu_{n}=\nu(x^{n})\) （对所有 \(n\in N\) ）。 f) 半双线性形式 \(b_{\mu}\) 是分离的，当且仅当对每个测度 \(\nu\in\mathcal{M}_{\mathrm{m}}(\mathbf{R})\) （满足 \(\mu_{n}=\nu(x^{n})\) ，对所有 \(n\in N\) ），\(\nu\) 的支撑都是无穷的。（注意不等式 (24) 的左端等于 \(b_{\mu}(p,p)\) ，其中 \(p=a_{0}+\cdots+a_{n}X^{n}\) 。）
\end{enumerate}

\begin{enumerate}
\def\labelenumi{\arabic{enumi})}
\setcounter{enumi}{21}
\tightlist
\item
  我们沿用上一习题的记号。
\end{enumerate}

设 \(\boldsymbol{\mu} = (\mu_n)_{n\in \mathbf{N}}\) 为满足条件 (24) 的实数序列。我们考虑测度 \(\nu \in \mathcal{M}_{\mathrm{m}}(\mathbf{R})\) （使 \(\mu_{n} = \nu (x^{n})\) ，对所有 \(n\in \mathbf{N}\) ）的唯一性问题。我们假设半双线性形式 \(b_{\boldsymbol{\mu}}\) 是分离的（参见上一习题的问题 \(f\) ））。

\begin{enumerate}
\def\labelenumi{\alph{enumi})}
\item
  若存在具有紧支撑的测度 \(\nu \in \mathcal{M}_{\mathrm{m}}(\mathbf{R})\) 使 \(\nu(x^{n}) = \mu_{n}\) （对所有 \(n \in \mathbf{N}\) ），则 \(\nu\) 是 \(\mathcal{M}_{\mathrm{m}}(\mathbf{R})\) 中使 \(\nu(x^{n}) = \mu_{n}\) （对所有 \(n \in \mathbf{N}\) ）成立的唯一测度。
\item
  设 \(\nu \in \mathcal{M}_{\mathrm{m}}(\mathbf{R})\) 满足 \(\nu (x^n) = \mu_n\) （对所有 \(n\in \mathbf{N}\) ）。对每个 \(p\in \mathbf{C}[\mathrm{X}]\) ，我们有 \(p\in \mathcal{L}^2 (\mathbf{R},\nu)\) 且 \(\| p\| = b_{\mu}(p,p)\) 。
\item
  从 C{[}X{]} 到 \(\mathcal{L}^{2}(\mathbf{R},\nu)\) 的典范单射通过连续性与过渡到商，定义了从 \(E_{\mu}\) 到 \(\mathrm{L}^{2}(\mathbf{R},\nu)\) 的一个等距。在下文中，我们把 \(E_{\mu}\) 与它的像等同起来，该像是 \(\mathrm{L}^{2}(\mathbf{R},\nu)\) 的一个闭子空间。
\end{enumerate}

\(d)\) 我们有 \(u_{\mu} \subset m_x\) ，其中 \(m_x\) 表示乘以 \(x\) 的算子（在 \(\mathrm{L}^2(\mathbf{R}, \nu)\) 中）。

\begin{enumerate}
\def\labelenumi{\alph{enumi})}
\setcounter{enumi}{4}
\item
  测度 \(\nu\) 是常函数 1 关于 \(m_x\) 的谱测度。
\item
  设 \(\nu \in \mathcal{M}_{\mathrm{m}}(\mathbf{R})\) 。我们有 \(\nu(x^n) = \mu_n\) （对所有 \(n \in \mathbf{N}\) ），当且仅当存在希尔伯特空间 F 与 F 上满足下列条件的自伴部分算子 \(v\) ：
\end{enumerate}

\begin{enumerate}
\def\labelenumi{(\roman{enumi})}
\item
  空间 \(E_{\mu}\) 是 F 的闭子空间，且从 \(E_{\mu}\) 到 F 的典范单射是等距的；
\item
  我们有 \(u_{\mu} \subset v\) ；
\item
  测度 \(\nu\) 是常函数 \(1 \in E_{\mu}\) 关于 v 的谱测度。
\end{enumerate}

\begin{enumerate}
\def\labelenumi{\alph{enumi})}
\setcounter{enumi}{6}
\item
  存在唯一的测度 \(\nu \in \mathcal{M}_{\mathrm{m}}(\mathbf{R})\) 使 \(\nu(x^{m}) = \mu_{n}\) （对所有 \(n \in \mathbf{N}\) ），当且仅当部分算子 \(u_{\boldsymbol{\mu}}\) 本质自伴。
\item
  设 \(\nu_{lg}\) 为 R 上的正测度，其支撑含于 \(R_{+}\) ，且它在 \(R_{+}^{*}\) 上的限制相对于勒贝格测度具有密度
\end{enumerate}

\[
t \mapsto \frac {1}{t} \exp (- \log (t) ^ {2} / 2)
\]

设 \(\widetilde{\nu}_{lg}\) 为 R 上的正测度，其支撑含于 \(R_{+}\) ，它在 \(R_{+}^{*}\) 上的限制以密度 \(t \mapsto 1 + \sin(2\pi \log(t))\) 相对于 \(\nu_{lg}\) 定义。

我们有 \(\nu_{\mathrm{lg}} \in \mathcal{M}_{\mathrm{m}}(\mathbf{R})\) 且 \(\widetilde{\nu}_{\mathrm{lg}} \in \mathcal{M}_{\mathrm{m}}(\mathbf{R})\) 。

\begin{enumerate}
\def\labelenumi{\roman{enumi})}
\tightlist
\item
  我们有 \(\nu_{\mathrm{lg}}(x^n) = \widetilde{\nu}_{\mathrm{lg}}(x^n)\) （对所有 \(n\in \mathbf{N}\) ）。
\end{enumerate}

\begin{enumerate}
\def\labelenumi{\arabic{enumi})}
\setcounter{enumi}{22}
\tightlist
\item
  我们沿用上一习题的记号；特别地，我们考虑实数序列 \(\boldsymbol{\mu} = (\mu_n)_{n\in \mathbf{N}}\) ，它满足条件 (24) 且使半双线性形式 \(b_{\mu}\) 分离。我们把 \(\mathbf{C}[\mathrm{X}]\) 与 \(\mathrm{E}_{\mu}\) 的一个稠密子空间等同起来。
\end{enumerate}

\begin{enumerate}
\def\labelenumi{\alph{enumi})}
\tightlist
\item
  存在唯一的多项式序列 \((p_{n})_{n\in\mathbf{N}}\) （属于 C{[}X{]}），使下列条件得到满足（“与 \(\mu\) 相关联的第一类正交多项式”）：
\end{enumerate}

\begin{enumerate}
\def\labelenumi{(\roman{enumi})}
\item
  族 \((p_n)\) 是 \(\mathrm{E}_{\mu}\) 的标准正交基；
\item
  多项式 \(p_n\) 的次数为 \(n\) （对所有 \(n \in \mathbf{N}\) ）；
\item
  \(\mathrm{X}^n\) 在 \(p_n\) 中的系数属于 \(\mathbf{R}_+^*\) 。
\end{enumerate}

此外，我们有 \(p_n \in \mathbf{R}[\mathbf{X}]\) （对所有 \(n \in \mathbf{N}\) ），且 \(p_0 = 1 / \sqrt{\mu_0}\) 。我们还记 \(p_{-1} = 0\) 。

\begin{enumerate}
\def\labelenumi{\alph{enumi})}
\setcounter{enumi}{1}
\tightlist
\item
  存在序列 \((a_{n})_{n\geqslant -1}\) （在 \(\mathbf{R}_{+}^{*}\) 中）与序列 \((b_{n})_{n\in \mathbf{N}}\) （在 \(\mathbf{R}\) 中），使 \(a_{-1} = 1\) 且
\end{enumerate}

\[
u _ {\pmb {\mu}} (p _ {n}) = a _ {n} p _ {n + 1} + b _ {n} p _ {n} + a _ {n - 1} p _ {n - 1}
\]

对所有 \(n \in \mathbf{N}\) 成立。这些序列是唯一的。

\begin{enumerate}
\def\labelenumi{\alph{enumi})}
\setcounter{enumi}{2}
\item
  对每个 \(n \geqslant 1\) ，\(X^{n}\) 在正交多项式 \(p_{n}\) 中的系数等于 \((a_{0} \cdots a_{n-1})^{-1}\) 。
\item
  设 \(\mathcal{F}\) 为由复数序列 \((z_{n})_{n\geqslant -1}\) （满足 \(z_{-1} = 0\) ）组成的复向量空间。我们用 \(\widetilde{v}_{\mu}\) 表示从 \(\mathcal{F}\) 到其自身的线性映射，它满足 \(\widetilde{v}_{\mu}(z_n) = (w_n)_{n\in \mathbf{N}}\) ，其中 \(w_{-1} = 0\) 且
\end{enumerate}

\[
w _ {n} = a _ {n} z _ {n + 1} + b _ {n} z _ {n} + a _ {n - 1} z _ {n - 1}
\]

对所有 \(n \in N\) 成立。

对每个 \(s \in \mathbf{C}\) ，序列 \(z_s = (p_n(s))_{n \geqslant -1}\) 是 \(\mathcal{F}\) 中唯一的满足 \((z_s)_0 = p_0(z)\) 且 \(\widetilde{v}_{\boldsymbol{\mu}}(z_s) = sz_s\) 的序列。

\begin{enumerate}
\def\labelenumi{\alph{enumi})}
\setcounter{enumi}{4}
\tightlist
\item
  设 \(z\) 与 \(w\) 为 \(\mathcal{F}\) 的元素。我们通过令下式成立来定义序列 \(\mathrm{W}(z,w) = (\mathrm{W}_n(z,w))_{n\in \mathbf{N}}\) ：
\end{enumerate}

\[
\mathrm{W} _ {n} (z, w) = a _ {n} (z _ {n + 1} w _ {n} - z _ {n} w _ {n + 1})
\]

对 \(n\in \mathbf{N}\) 成立。我们有

\[
\mathrm{W} _ {n} (z, w) = \sum_ {j = 0} ^ {n} (\widetilde {v} _ {\boldsymbol {\mu}} (z) _ {j} w _ {j} - z _ {j} \widetilde {v} _ {\boldsymbol {\mu}} (w) _ {j})
\]

对所有 \(n \in N\) 成立。

\(f)\) 设 \(\mathrm{P} \subset \ell^2(\mathbf{N})\) 为具有有限支撑的序列组成的空间。对每个 \(z \in \mathrm{P}\) ，序列 \(\widetilde{v}_{\boldsymbol{\mu}}(z)\) 属于 \(\ell^2(\mathbf{N})\) 。

我们用 \(v_{\mu}\) 表示 \(\ell^2(\mathbf{N})\) 上以 P 为定义域、由 \(z \mapsto \widetilde{v}_{\mu}(z)\) 给出的部分算子。它是对称的。

\begin{enumerate}
\def\labelenumi{\alph{enumi})}
\setcounter{enumi}{6}
\tightlist
\item
  \(v_{\mu}\) 的伴随是部分算子，其定义域为
\end{enumerate}

\[
\mathrm{dom} (v _ {\boldsymbol {\mu}} ^ {*}) = \{z \in \ell^ {2} (\mathbf {N}) | \widetilde {v} _ {\boldsymbol {\mu}} (z) \in \ell^ {2} (\mathbf {N}) \}
\]

且满足 \(v_{\mu}^{*}(z) = \widetilde{v}_{\mu}(z)\) （对 \(z \in \mathrm{dom}(v_{\mu}^{*})\) ）。

\(h)\) 对所有 \(z\) 与 \(w\) （属于 \(\mathrm{dom}(v_{\mu}^{*})\) ），\(\mathrm{W}_{n}(\overline{z}, w)\) 当 \(n \to +\infty\) 时的极限存在，且等于 \(\langle w | v_{\mu}^{*}(z) \rangle - \langle v_{\mu}^{*}(w) | z \rangle\) 。

\begin{enumerate}
\def\labelenumi{\roman{enumi})}
\tightlist
\item
  设 \(s \in C\) 。\(v_{\boldsymbol{\mu}}^{*} - s \cdot 1_{\ell^{2}(\mathbf{N})}\) 的核维数 \(\leqslant 1\) ；其维数为 1 当且仅当 \(z_{s} \in \ell^{2}(\mathbf{N})\) ，且此时它由 \(z_{s}\) 生成。
\end{enumerate}

\begin{enumerate}
\def\labelenumi{\alph{enumi})}
\setcounter{enumi}{9}
\tightlist
\item
  部分算子 \(v_{\mu}\) （相应地 \(u_{\mu}\) ）的闭包的亏指标为 (0, 0) 或 (1, 1)；第一种情形成立当且仅当存在 \(s \in \mathbf{C} - \mathbf{R}\) 使 \(z_s \notin \ell^2(\mathbf{N})\) 。（对 \(u_{\mu}\) 的情形，注意映射 \(X^n \mapsto e_n\) （其中 \((e_n)\) 是 \(\ell^2(\mathbf{N})\) 的典范基）诱导出 \(E_{\mu}\) 到 \(\ell^2(\mathbf{N})\) 的一个等距同构，使我们可以把 \(u_{\mu}\) 与 \(v_{\mu}\) 等同起来。）
\end{enumerate}

\(k)\) 设 \(s \in \mathbf{C}\) ，\(f \in \mathrm{Ker}(u_{\mu}^{*} - s \cdot 1_{\mathrm{E}_{\mu}})\) 。若 \(\langle f | 1 \rangle = 0\) ，则 \(f = 0\) 。

\begin{enumerate}
\def\labelenumi{\alph{enumi})}
\setcounter{enumi}{11}
\item
  设 \(u_1\) 与 \(u_2\) 为 \(u_\mu\) 在 \(\mathrm{E}_\mu\) 上的两个自伴扩张。我们有 \(u_1 = u_2\) ，当且仅当存在 \(s \in \mathbf{C} - \mathbf{R}\) 使 \(\langle \mathrm{R}(u_1, s)1|1\rangle = \langle \mathrm{R}(u_2, s)1|1\rangle\) 。（假设这样一个等式成立；令 \(f_1 = \mathrm{R}(u_1, s)1\) ，并证明 \(f_1\) 不属于 \(\overline{u}_\mu\) 的定义域；由此推出 \(u_1\) 与 \(u_1\) 到 \(\mathrm{dom}(\overline{u}_\mu) \oplus \mathbf{C}f_1\) 的限制的闭包重合。对 \(u_2\) 与 \(f_2 = \mathrm{R}(u_2, s)1\) 作同样的推理，然后证明 \(f_1 = f_2\) ，从而得出 \(u_1 = u_2\) 。） m) 假设 \(u_\mu\) 不是本质自伴的。则至少存在两个测度 \(\nu \in \mathcal{M}_{\mathrm{m}}(\mathbf{R})\) （满足 \(\nu(x^n) = \mu_n\) ，对所有 \(n \in \mathbf{N}\) ）。 n) 以下设 \(u_\mu\) 本质自伴。对每个 \(s \in \mathbf{C} - \mathbf{R}\) ，存在多项式序列 \((r_{s,m})_{m \in \mathbf{N}}\) ，使 \((x - s)r_{s,m}\) 在 \(\mathrm{E}_\mu\) 中当 \(m \to +\infty\) 时收敛于常函数 1。
\item
  设 \(\nu\in\mathcal{M}_{\mathrm{m}}(\mathbf{R})\) 为满足 \(\nu(x^{n})=\mu_{n}\) （对所有 \(n\in\mathbf{N}\) ）的测度。对每个 \(m\in\mathbf{N}\) ，不等式
\end{enumerate}

\[
\left| \int_ {\mathbf {R}} \frac {d \nu (t)}{t - s} - \int_ {\mathbf {R}} r _ {s, m} (t) d \nu (t) \right| \leqslant \frac {\mu_ {0} ^ {2}}{| \mathcal {I} (s) ^ {2} |} \| 1 - (x - s) r _ {s, m} \| _ {\mathrm{E} _ {\mu}}
\]

成立。

\begin{enumerate}
\def\labelenumi{\alph{enumi})}
\setcounter{enumi}{15}
\tightlist
\item
  测度 \(\nu\) 是唯一的。
\end{enumerate}

\begin{enumerate}
\def\labelenumi{\arabic{enumi})}
\setcounter{enumi}{23}
\tightlist
\item
  我们沿用前面诸习题的记号；特别地，我们考虑实数序列 \(\boldsymbol{\mu} = (\mu_n)_{n\in \mathbf{N}}\) ，它满足条件 (24) 且使半双线性形式 \(b_{\boldsymbol{\mu}}\) 分离。
\end{enumerate}

\begin{enumerate}
\def\labelenumi{\alph{enumi})}
\tightlist
\item
  存在唯一的具有实系数的多项式序列 \((q_{n})_{n\in\mathbf{N}}\) ，使对每个 \(s\in C\) ，公式
\end{enumerate}

\[
q _ {n} (s) = b _ {\pmb {\mu}} \Big (1, \frac {p _ {n} - p _ {n} (s)}{\mathrm{X} - s} \Big)
\]

成立（“第二类正交多项式”）。

\begin{enumerate}
\def\labelenumi{\alph{enumi})}
\setcounter{enumi}{1}
\tightlist
\item
  设 \(\nu \in \mathcal{M}_{\mathrm{m}}(\mathbf{R})\) 为 \(\mathbf{R}\) 上满足 \(\nu(x^n) = \mu_n\) （对所有 \(n \in \mathbf{N}\) ）的测度。对每个 \(s \in \mathbf{C} - \mathbf{R}\) ，我们有
\end{enumerate}

\[
\int_ {\mathbf {R}} \frac {p _ {n} (t)}{t - s} d \nu (t) = q _ {n} (s) + p _ {n} (s) \int_ {\mathbf {R}} \frac {d \nu (t)}{t - s}.
\]

\begin{enumerate}
\def\labelenumi{\alph{enumi})}
\setcounter{enumi}{2}
\item
  设 \(s \in C-R\) 。序列 \((p_{n}(s))_{n \in \mathbf{N}}\) 属于 \(\ell^{2}(\mathbf{N})\) ，当且仅当序列 \((q_{n}(s))_{n \in \mathbf{N}}\) 属于 \(\ell^{2}(\mathbf{N})\) 。（令 \(\sigma(s) = \int(t-s)^{-1} d\nu(t)\) ；利用上一问题，证明序列 \((q_{n}(s) + \sigma(s)p_{n}(s))_{n \in \mathbf{N}}\) 属于 \(\ell^{2}(\mathbf{N})\) 。）
\item
  对每个 \(n \in \mathbf{N}\) ，公式
\end{enumerate}

\[
a _ {n} (p _ {n + 1} (\mathrm{X}) p _ {n} (\mathrm{Y}) - p _ {n} (\mathrm{X}) p _ {n + 1} (\mathrm{Y})) = (\mathrm{X} - \mathrm{Y}) \sum_ {j = 0} ^ {n} p _ {j} (\mathrm{X}) p _ {j} (\mathrm{Y})
\]

在 \(\mathbf{C}[\mathrm{X},\mathrm{Y}]\) 中成立。

\begin{enumerate}
\def\labelenumi{\alph{enumi})}
\setcounter{enumi}{4}
\tightlist
\item
  对每个 \(n \in \mathbf{N}\) ，我们有关系式
\end{enumerate}

\[
p _ {n} q _ {n + 1} - p _ {n + 1} q _ {n} = \frac {1}{a _ {n}}.
\]

\begin{enumerate}
\def\labelenumi{\alph{enumi})}
\setcounter{enumi}{5}
\tightlist
\item
  若通项为 \(a_{n}^{-1}\) 的级数发散，则测度 \(\nu\) 是满足 \(\nu(x^{n}) = \mu_{n}\) （对所有 \(n \in N\) ）的唯一测度。
\end{enumerate}

\begin{enumerate}
\def\labelenumi{\arabic{enumi})}
\setcounter{enumi}{24}
\tightlist
\item
  我们沿用前面诸习题的记号；特别地，我们考虑实数序列 \(\boldsymbol{\mu} = (\mu_{n})_{n \in \mathbf{N}}\) ，它满足条件 (24) 且使半双线性形式 \(b_{\mu}\) 分离。
\end{enumerate}

假设 \(\mu_{0}=1\) ，且卡尔曼条件成立，即通项为 \(\mu_{2n}^{-1/(2n)}\) （\(n\geqslant1\) ）的级数发散。 a) 我们有

\[
\mu_ {2 n} ^ {- 1 / (2 n)} \leqslant (a _ {0} \dots a _ {n - 1}) ^ {- 1 / n} \leqslant \frac {e}{n ^ {2}} \sum_ {j = 1} ^ {n} \frac {j}{a _ {j - 1}}
\]

对每个整数 \(n \geqslant 1\) 成立。（为建立第二个不等式，可以证明并利用不等式 \((n/e)^{n} \leqslant n!\) （\(n \geqslant 1\) ）。）

\begin{enumerate}
\def\labelenumi{\alph{enumi})}
\setcounter{enumi}{1}
\item
  通项为 \(a_{n}^{-1}\) （\(n \geqslant 1\) ）的级数发散。
\item
  存在唯一的测度 \(\nu \in \mathcal{M}_{\mathrm{m}}(\mathbf{R})\) ，满足 \(\nu(x^n) = \mu_n\) （对所有 \(n \in \mathbf{N}\) ）。
\end{enumerate}

\begin{enumerate}
\def\labelenumi{\arabic{enumi})}
\setcounter{enumi}{25}
\tightlist
\item
  我们用 \(x\) 表示 \(\mathbf{R}_{+}\) 的恒等函数。
\end{enumerate}

\begin{enumerate}
\def\labelenumi{\alph{enumi})}
\tightlist
\item
  设 \(\nu\) 为 \(R_{+}\) 上的正测度，使 \(x^{n}\) 是 \(\nu\) -可积的（对每个整数 \(n \geqslant 0\) ）。记 \(\mu_{n} = \nu(x^{n})\) 。对每个整数 \(n \geqslant 0\) 与所有复数 \(a_{0}, \ldots, a_{n}\) ，我们有
\end{enumerate}

\[
\sum_ {i = 0} ^ {n} \sum_ {j = 0} ^ {n} \overline {{a}} _ {i} a _ {j} \mu_ {i + j} \geqslant 0
\]

\[
\sum_ {i = 0} ^ {n} \sum_ {j = 0} ^ {n} \overline {{a}} _ {i} a _ {j} \mu_ {i + j + 1} \geqslant 0.
\]

\begin{enumerate}
\def\labelenumi{\alph{enumi})}
\setcounter{enumi}{1}
\tightlist
\item
  若 \((\mu_{n})_{n\geqslant0}\) 是满足上一问题条件的实数序列，则存在正测度 \(\nu\) （在 \(R_{+}\) 上），使 \(\mu_{n}=\nu(x^{n})\) （对所有 \(n\in N\) ）。（“斯蒂尔切斯矩问题”；利用习题 21 的方法，并借助弗里德里希斯扩张。）
\end{enumerate}

\begin{enumerate}
\def\labelenumi{\arabic{enumi})}
\setcounter{enumi}{26}
\tightlist
\item
  我们记 \(M = X^{2}Y^{2}(X^{2} + Y^{2} - 3) \in \mathbf{R}[X, Y]\) （“莫茨金多项式”）。
\end{enumerate}

\begin{enumerate}
\def\labelenumi{\alph{enumi})}
\item
  我们有 M(x,y) ≥ 0 （对所有 x, y ∈ R）。
\item
  不存在有限族 \((p_{i})_{i\in I}\) （在 R{[}X,Y{]} 中）使 \(M=\sum p_{i}^{2}\) 。（若这样一个族存在，证明每个 \(p_{i}\) 必形如 \(a+bXY\) ，其中 \(a\in R\) 且 \(b\in R[X,Y]\) 的次数至多为 1；然后计算 \((\mathrm{XY})^{2}\) 的系数。）
\item
  设 \(\tau\colon\mathbf{N}^{2}\to\mathbf{N}-\{0\}\) 为由下式定义的映射：
\end{enumerate}

\[
\tau (0, 0) = 1, \quad \tau (1, 2) = 2, \quad \tau (2, 1) = 3, \quad \tau (1, 1) = 4, \quad \tau (1, 0) = 5,
\]

\[
\tau (0, 1) = 6, \quad \tau (2, 0) = 7, \quad \tau (0, 2) = 8, \quad \tau (3, 0) = 9, \quad \tau (0, 3) = 1 0
\]

以及

\[
\tau (i, j) = j + 1 + (i + j) (i + j + 1) / 2 \text {   if   } i + j \geqslant 4.
\]

映射 \(\tau\) 是一个双射。

\begin{enumerate}
\def\labelenumi{\alph{enumi})}
\setcounter{enumi}{3}
\tightlist
\item
  设 \((\nu_{n})_{n\geqslant 1}\) 为由下式定义的序列：
\end{enumerate}

\[
\nu_ {i} = 1 \text {if} i \in \{1, 2, 3 \}, \quad \nu_ {4} = 4, \quad \nu_ {n} = (n!) ^ {(n + 1)!} \text {if} n \geqslant 5.
\]

对 \(n\) 与 \(m\) （属于 \(\mathbf{N}\) ），我们令 \(\mu_{n,m} = 0\) （若 \(n\) 或 \(m\) 为奇数），否则令 \(\mu_{n,m} = \nu_{\tau(n/2,m/2)}\) 。线性映射 \(\alpha\) （从 \(\mathbf{C}[\mathrm{X},\mathrm{Y}]\) 到 \(\mathbf{C}\) ，满足 \(\alpha(\mathrm{X}^n\mathrm{Y}^m) = \mu_{n,m}\) ）满足 \(\alpha(p^2) \geqslant 0\) （对所有 \(p \in \mathbf{R}[\mathrm{X},\mathrm{Y}]\) ）。（只需验证：对每个整数 \(k \in \mathbf{N}\) ，矩阵 \((a_{m,m})_{1 \leqslant n,m \leqslant k}\) 的行列式为正，其中 \(a_{n,m}\) 由 \(a_{\tau(i,j),\tau(i',j')} = \alpha(\mathrm{X}^{i+i'}\mathrm{Y}^{j+j'})\) （对所有 \((i,j,i',j')\) ）定义；然后对 \(k\) 归纳。）

\begin{enumerate}
\def\labelenumi{\alph{enumi})}
\setcounter{enumi}{4}
\tightlist
\item
  不存在正测度 \(\nu\) （在 \(R^{2}\) 上），使得 \((x,y)\mapsto x^{n}y^{m}\) 关于 \(\mu\) 可积（对所有 \((n,m)\in\mathbf{N}^{2}\) ），且使得
\end{enumerate}

\[
\int_ {\mathbf {R} ^ {2}} x ^ {n} y ^ {m} d \nu (x, y) = \mu_ {n, m}
\]

对所有 \((n,m)\in\mathbf{N}^{2}\) 成立。

\begin{enumerate}
\def\labelenumi{\alph{enumi})}
\setcounter{enumi}{5}
\tightlist
\item
  存在希尔伯特空间 E 与 E 上一个稠定部分算子 u，满足下列条件：
\end{enumerate}

\begin{enumerate}
\def\labelenumi{(\roman{enumi})}
\item
  我们有 \(\operatorname{dom}(u) \subset \operatorname{dom}(u^*)\) ；
\item
  我们有 \(\|u(x)\|=\|u^{*}(x)\|\) （对所有 \(x\in\mathrm{dom}(u)\) ）；
\item
  不存在包含 E 为子空间的希尔伯特空间 F（使 E 到 F 的典范单射是等距的）以及 F 上的正规部分算子 v 使 \(u \subset v\) 。
\end{enumerate}

把这一结果与 IV 第 352 页习题 24 相比较。

\begin{enumerate}
\def\labelenumi{\arabic{enumi})}
\setcounter{enumi}{27}
\tightlist
\item
  对每个整数 \(k \geqslant 0\) ，我们用 \(b_k\) 表示集合 \(\{1, \ldots, k\}\) 上等价关系的个数。
\end{enumerate}

对每个 \(\lambda > 0\) ，我们用 \(\varpi_{\lambda}\) 表示参数为 \(\lambda\) 的泊松测度（在 \(\mathbf{R}\) 上）（参见 II 第 282 页习题 41），即总质量为 1 、支撑在 \(\mathbf{N}\) 上且满足 \(\varpi_{\lambda}(j) = e^{-\lambda} \lambda^{j} / j!\) （对所有 \(j \in \mathbf{N}\) ）的测度。

\begin{enumerate}
\def\labelenumi{\alph{enumi})}
\tightlist
\item
  关系式
\end{enumerate}

\[
b _ {k + 1} = \sum_ {j = 0} ^ {k} {\binom {k} {j}} b _ {j}
\]

对每个整数 \(k \in N\) 成立。

\begin{enumerate}
\def\labelenumi{\alph{enumi})}
\setcounter{enumi}{1}
\tightlist
\item
  对任意函数 \(f: N \to C\) （使函数 \(x \mapsto xf(x)\) 是 \(\varpi\) -可积的），我们有
\end{enumerate}

\[
\int_ {\mathbf {R}} x f (x) d \varpi_ {\lambda} (x) = \lambda \int_ {\mathbf {R}} f (x + 1) d \varpi_ {\lambda} (x)
\]

（“陈氏公式”）。

此外，若 \(\mu\) 是 R 上满足这一性质的有界测度，则存在 \(c \in R\) 使 \(\mu = c\varpi_{\lambda}\) 。

\begin{enumerate}
\def\labelenumi{\alph{enumi})}
\setcounter{enumi}{2}
\tightlist
\item
  对每个 \(k \in \mathbf{N}\) ，函数 \(x \mapsto x^k\) 是 \(\varpi_1\) -可积的，且我们有
\end{enumerate}

\[
\varpi_ {1} (x \mapsto x ^ {k}) = b _ {k}.
\]

此外，测度 \(\varpi_{1}\) 是 R 上满足这些等式的唯一正测度。

\begin{enumerate}
\def\labelenumi{\alph{enumi})}
\setcounter{enumi}{3}
\tightlist
\item
  对每个整数 \(k \geqslant 0\) ，我们有
\end{enumerate}

\[
b _ {k} = \frac {1}{e} \sum_ {j \geqslant 0} \frac {j ^ {n}}{j !}
\]

（“多宾斯基公式”）以及

\[
b _ {k} = \frac {2 k !}{e \pi} \mathcal {I} \left(\int_ {0} ^ {\pi} e ^ {e ^ {e ^ {i t}}} \sin (k t) d t\right)
\]

（“切萨罗公式”）。

\begin{enumerate}
\def\labelenumi{\alph{enumi})}
\setcounter{enumi}{4}
\tightlist
\item
  对每个整数 \(n \geqslant 1\) ，我们用 \(\mu_{n}\) 表示 R 上的正测度，它是对称群 \(S_{n}\) 的规范化哈尔测度在映射 \(f_{n}\) （从 \(S_{n}\) 到 R）下的像，该映射把 \(\sigma \in S_{n}\) 映为 \(\sigma\) 的不动点个数。设 \(k \in N\) 。对每个整数 \(n \geqslant k\) ，我们有
\end{enumerate}

\[
\mu_ {n} (x \mapsto x ^ {k}) = b _ {k}.
\]

\begin{enumerate}
\def\labelenumi{\alph{enumi})}
\setcounter{enumi}{5}
\tightlist
\item
  序列 \((\mu_{n})\) 当 \(n \to +\infty\) 时弱收敛于泊松测度 \(\varpi_{1}\) 。
\end{enumerate}

\begin{enumerate}
\def\labelenumi{\arabic{enumi})}
\setcounter{enumi}{28}
\tightlist
\item
  设 E 为复希尔伯特空间。我们用 \(\mathcal{T}_b\) （相应地 \(\mathcal{T}_s\) 、 \(\mathcal{T}_f\) ）表示 \(\mathcal{L}(\mathrm{E})\) 的巴拿赫空间拓扑（相应地逐点收敛拓扑、由半范数 \(u \mapsto \langle x \mid u(y) \rangle\) （对 \((x, y) \in \mathrm{E} \times \mathrm{E}\) ）定义的局部凸拓扑）。
\end{enumerate}

\begin{enumerate}
\def\labelenumi{\alph{enumi})}
\item
  设 \(\lambda_0\in \mathbf{C} - \mathbf{R}\) 。 \(\mathcal{T}_b\) -预解拓扑（相应地 \(\mathcal{T}_s\) -预解拓扑）是 \(\mathcal{A}(\mathrm{E})\) 上使映射 \(u\mapsto \mathrm{R}(u,\lambda_0)\) （从 \(\mathcal{A}(\mathrm{E})\) 到 \(\mathcal{L}(\mathrm{E})\) （赋有拓扑 \(\mathcal{T}_b\) （相应地 \(\mathcal{T}_s\) ）））连续的最粗拓扑。
\item
  \(T_{f}\) -预解拓扑与 \(T_{s}\) -预解拓扑重合。
\item
  设 \(X \subset \mathcal{L}(E)\) 为 E 的埃尔米特自同态组成的有界集。由 \(T_{b}\) -预解拓扑（相应地 \(T_{s}\) -预解拓扑）在 X 上诱导的拓扑，与由 \(T_{b}\) （相应地 \(T_{s}\) ）诱导的拓扑重合。
\item
  a) 中的性质对 \(T_{f}\) -预解拓扑不成立。
\item
  给出一个序列 \((u_{n})\) （由部分自伴算子组成，且对 \(T_{b}\) -预解拓扑收敛）以及函数 \(f \in \mathcal{C}_{b}(\mathbf{R})\) 的例子，使 \(f(u_{n})\) 不收敛于 \(f(u)\) （在 \(\mathcal{L}(\mathrm{E})\) 中）。
\end{enumerate}

\(f\) 设 \((u_n)\) 为 E 上部分自伴算子的序列，\(u\) 为 E 上的一个自伴部分算子。若存在 \(\mathrm{E}_u\) 的稠密子空间 F（包含于 \(\mathrm{dom}(u_n)\) ，对所有 \(n\) ），使 \(u_n(x)\) 收敛于 \(u(x)\) （对所有 \(x \in \mathrm{F}\) ），则 \((u_n)\) 收敛于 \(u\) （对 \(\mathcal{T}_s\) -预解拓扑）。

\begin{enumerate}
\def\labelenumi{\arabic{enumi})}
\setcounter{enumi}{29}
\tightlist
\item
  设 E 为复希尔伯特空间。用 \(T_{b}\) 表示 \(\mathcal{L}(\mathrm{E})\) 上有界收敛拓扑。
\end{enumerate}

\begin{enumerate}
\def\labelenumi{\alph{enumi})}
\tightlist
\item
  设 \(u\) 为 E 上的自伴部分算子。对每个 \(\lambda \in \mathbf{C} - \mathbf{R}\) （使 \(\lambda\) 的虚部为负），我们有
\end{enumerate}

\[
\mathrm{R} (u, \lambda) = i \int_ {0} ^ {\infty} e ^ {- i t \lambda} e ^ {i t u} d t.
\]

\begin{enumerate}
\def\labelenumi{\alph{enumi})}
\setcounter{enumi}{1}
\tightlist
\item
  \(T_{b}\) -预解拓扑是 \(\mathcal{A}(E)\) 上使映射 \(u \mapsto e^{itu}\) （从 \(\mathcal{A}(E)\) 到 \(\mathcal{L}(E)\) ）对每个 \(t \in R\) 都连续的最粗拓扑。
\end{enumerate}

\begin{enumerate}
\def\labelenumi{\arabic{enumi})}
\setcounter{enumi}{30}
\tightlist
\item
  设 E 为复希尔伯特空间。用 \(T_{s}\) 表示 \(\mathcal{L}(\mathrm{E})\) 的逐点收敛拓扑。用 H 表示虚部严格为正的复数组成的集。
\end{enumerate}

设 \((u_{n})\) 为 \(\mathcal{A}(\mathrm{E})\) 中的序列。我们假设

\begin{enumerate}
\def\labelenumi{(\roman{enumi})}
\item
  存在 \(\lambda_{+} \in \mathbf{H}\) 使 \(\mathrm{R}(u_n, \lambda_{+})\) 在 \(\mathcal{L}(\mathrm{E})\) （赋有拓扑 \(\mathcal{T}_s\) ）中收敛；我们用 \(v_{+}\) 表示其极限；
\item
  存在 \(\lambda_{-} \in \overline{\mathbf{H}}\) 使 \(\mathrm{R}(u_n, \lambda_{-})\) 在 \(\mathcal{L}(\mathrm{E})\) （赋有拓扑 \(\mathcal{T}_s\) ）中收敛；
\item
  自同态 \(v_{+}\) 的像 F 在 E 中稠密。
\end{enumerate}

设 D 为 C 中由满足下式的 \(\lambda \in H\) 组成的开圆盘

\[
| \lambda - \lambda_ {+} | <   \frac {1}{| \mathcal {I} (\lambda_ {+}) |}.
\]

\begin{enumerate}
\def\labelenumi{\alph{enumi})}
\tightlist
\item
  对每个 \(\lambda \in D\) ，级数
\end{enumerate}

\[
r _ {\lambda} = \sum_ {n = 0} ^ {\infty} (\lambda_ {+} - \lambda) v _ {+} ^ {n + 1}
\]

在 \(\mathcal{L}(\mathrm{E})\) 中收敛。此外，对每个 \(\lambda \in \mathrm{D}\) ，序列 \((\mathrm{R}(u_n,\lambda))_n\) 收敛于 \(r_{\lambda}\) （在 \(\mathcal{L}(\mathrm{E})\) 中，赋有拓扑 \(\mathcal{T}_s\) ）。

\begin{enumerate}
\def\labelenumi{\alph{enumi})}
\setcounter{enumi}{1}
\item
  存在唯一的从 H 到 \(\mathcal{L}(\mathrm{E})\) 的全纯映射，它在 D 上与映射 \(\lambda\mapsto r_{\lambda}\) 重合。我们仍把它记作 \(\lambda\mapsto r_{\lambda}\) 。
\item
  存在唯一的从 \(\mathbf{C} - \mathbf{R}\) 到 \(\mathcal{L}(\mathrm{E})\) 的全纯映射，它在 \(\mathbf{H}\) 上与映射 \(\lambda \mapsto r_{\lambda}\) 重合，且满足 \(r_{\lambda}^{*} = r_{\overline{\lambda}}\) （对所有 \(\lambda \in \mathbf{C} - \mathbf{R}\) ）。它满足 \(r_{\lambda_{-}} = v_{-}\) ，且对所有 \(\lambda \in \mathbf{C} - \mathbf{R}\) ，序列 \((\mathrm{R}(u_n,\lambda))\) 收敛于 \(r_{\lambda}\) （在 \(\mathcal{L}(\mathrm{E})\) 中，赋有拓扑 \(\mathcal{T}_s\) ）。
\item
  对所有 \(\lambda\) 与 \(\mu\) （属于 \(\mathbf{C} - \mathbf{R}\) ），我们有 \(r_{\lambda} - r_{\mu} = (\mu - \lambda) = r_{\mu}r_{\lambda}\) 以及 \(r_{\lambda}r_{\mu} = r_{\mu}r_{\lambda}\) 。\\
\item
  \(r_{\lambda}\) 的像对所有 \(\lambda \in \mathbf{C} - \mathbf{R}\) 都等于 F。
\end{enumerate}

\(f)\) 对每个 \(\lambda\in\mathbf{C}\text{—}\mathbf{R}\) ，自同态 \(r_{\lambda}\) 是单射；以 F 为定义域的部分算子 \(\lambda1_{\mathrm{E}}-r_{\lambda}^{-1}\) 不依赖于 \(\lambda\in\mathbf{C}\text{—}\mathbf{R}\) 的选取；我们把它记作 \(u\) 。 \(g)\) 部分算子 \(u\) 是自伴的；对所有 \(\lambda\in\mathbf{C}\text{—}\mathbf{R}\) ，我们有 R( \(u,\lambda\) ) = \(r_{\lambda}\) 。

\begin{enumerate}
\def\labelenumi{\alph{enumi})}
\setcounter{enumi}{7}
\tightlist
\item
  序列 \((u_{n})\) 对 \(T_{s}\) -预解拓扑收敛于 u。（“加藤–特罗特定理”）
\end{enumerate}

\begin{enumerate}
\def\labelenumi{\arabic{enumi})}
\setcounter{enumi}{31}
\tightlist
\item
  设 E 为复希尔伯特空间。设 u 与 v 为 E 上的部分自伴算子，使 \(u + v\) 自伴。
\end{enumerate}

\begin{enumerate}
\def\labelenumi{\alph{enumi})}
\item
  对 \(t \in \mathbf{R}^*\) ，令 \(v(t) = t^{-1}(e^{itu}e^{itv} - e^{it(u + v)}) \in \mathcal{L}(\mathrm{E})\) ，并用 \(w(t)\) 表示 \(v(t)\) 到希尔伯特空间 \(\mathrm{E}_{u + v}\) 的限制。用 B 表示连续线性映射 \(w(t)\) 组成的集。集 B 在巴拿赫空间 \(\mathcal{L}(\mathrm{E}_{u + v},\mathrm{E})\) 中有界。（利用巴拿赫–斯坦豪斯定理。）
\item
  我们有 \(w(t) \to 0\) （当 \(t \to 0\) ，在 \(R^{*}\) 中），收敛是在赋有紧收敛拓扑的空间 \(\mathcal{L}(\mathrm{E}_{u+v},\mathrm{E})\) 中取的。
\item
  设 \(x\in \mathrm{E}_{u + v}\) 。我们有
\end{enumerate}

\[
\lim _ {n \to + \infty} (e ^ {i t u / n} e ^ {i t v / n}) ^ {n} x = e ^ {i t (u + v)} x.
\]

对 \(n\in \mathbf{N}\) ，记

\[
\begin{array}{l} (e ^ {i t u / n} e ^ {i t v / n}) ^ {n} - e ^ {i t (u + v)} = \\ \sum_ {i = 0} ^ {n - 1} (e ^ {i t u / n} e ^ {i t v / n}) ^ {k} \Big (e ^ {i t u / n} e ^ {i t v / n} - e ^ {i t (u + v) / n} \Big) (e ^ {i t (u + v) / n}) ^ {n - 1 - i} \end{array}
\]

并把上一问题的结果应用于映射 \(t \mapsto e^{it(u+v)}x\) （从 \([-1,1]\) 到 \(E_{u+v}\) ）的紧像。

\begin{enumerate}
\def\labelenumi{\arabic{enumi})}
\setcounter{enumi}{32}
\tightlist
\item
  设 \(n \in \mathbf{N}\) 。我们用 \(\Delta\) 表示 \(\mathrm{L}^2(\mathbf{R}^n)\) 上的拉普拉斯算子。
\end{enumerate}

\begin{enumerate}
\def\labelenumi{\alph{enumi})}
\tightlist
\item
  对每个 \(t \in R^{*}\) ，自同态 \(e^{it\Delta}\) （属于 \(\mathrm{L}^{2}(\mathbf{R}^{n})\) ）由下式给出
\end{enumerate}

\[
e ^ {i t \Delta} f (x) = \frac {1}{(2 i t) ^ {n / 2}} \int_ {\mathbf {R} ^ {n}} f (y) e ^ {i | x - y | ^ {2} / (4 t)} d y
\]

若 \(f\in \mathcal{L}^1 (\mathbf{R}^n)\cap \mathcal{L}^2 (\mathbf{R}^n)\) （利用傅里叶变换。）

\begin{enumerate}
\def\labelenumi{\alph{enumi})}
\setcounter{enumi}{1}
\tightlist
\item
  设 C 为 \(\mathrm{L}^{2}(\mathbf{R}^{n})\) 上的酉算子，它由映射 \(f \mapsto \mathrm{C}(f)\) （对 \(f \in \mathcal{L}^{2}(\mathbf{R}^{n})\) 定义）通过过渡到商得到
\end{enumerate}

\[
\mathrm{C} (f) (x) = \frac {1}{(2 i t) ^ {n / 2}} e ^ {\frac {i x ^ {2}}{4 t}} \mathcal {F} f \Big (\frac {x}{2 t} \Big)
\]

对 \(x\in \mathbf{R}^n\)

我们有 \(e^{it\Delta} \to C\) （当 \(t \to +\infty\) ），收敛是在赋有单收敛拓扑的空间 \(\mathcal{L}(\mathrm{L}^{2}(\mathbf{R}^{n}))\) 中取的。

\begin{enumerate}
\def\labelenumi{\alph{enumi})}
\setcounter{enumi}{2}
\tightlist
\item
  假设 \(1 \leqslant n \leqslant 3\) 。设 \(g \in \mathcal{L}^{\infty}(\mathbf{R}^{n}) + \mathcal{L}^{2}(\mathbf{R}^{n})\) 为实值函数。用 u 表示部分算子 \(u = \Delta + m_{g}\) ；它在 \(\mathrm{L}^{2}(\mathbf{R}^{n})\) 上是自伴的（习题 14）。对 \(f \in \mathcal{L}^{1}(\mathbf{R}^{n}) \cap \mathcal{L}^{2}(\mathbf{R}^{n})\) 与 \(t \in R^{*}\) ，我们有
\end{enumerate}

\[
e ^ {i t u} f (x) = \lim _ {k \rightarrow + \infty} \frac {1}{(2 i t / k) ^ {k n / 2}} \int_ {\mathbf {R} ^ {n k}} f (y _ {k}) \exp \left(\frac {i t}{k} \sum_ {i = 1} ^ {k} \left(\frac {| y _ {i} - y _ {i - 1} | ^ {2}}{4 (t / k) ^ {2}} - g (y _ {i})\right) d y \right.
\]

其中我们记 \(y = (y_{1}, \ldots, y_{k})\) ，\(y_{i} \in R^{n}\) 。

\begin{enumerate}
\def\labelenumi{\alph{enumi})}
\setcounter{enumi}{3}
\tightlist
\item
  对上一公式的解释。
\end{enumerate}

\begin{enumerate}
\def\labelenumi{\arabic{enumi})}
\setcounter{enumi}{33}
\tightlist
\item
  设 E 为复希尔伯特空间。对 E 的所有自同态 \(u\) 与 \(v\) ，我们记 \([u,v] = uv - vu\) ，并定义 \(\mathrm{E}_{u,v} = \overline{\mathrm{Im}(u)} \cap \mathrm{Ker}(v)\) 。
\end{enumerate}

设 \(p\) 与 \(q\) 为 E 的正交投影算子。我们令 \(a = p - q\) 且 \(b = 1_{\mathrm{E}} - p - q\) 。 a) 若存在 E 的酉自同态 \(u\) 使 \(upu^{-1} = q\) 且 \(uqu^{-1} = p\) ，则 \(\mathrm{E}_{p,q}\) 与 \(\mathrm{E}_{1 - p,1 - q}\) 的希尔伯特维数相等。

\begin{enumerate}
\def\labelenumi{\alph{enumi})}
\setcounter{enumi}{1}
\tightlist
\item
  我们有关系式
\end{enumerate}

\[
a ^ {2} + b ^ {2} = 1 _ {\mathrm{E}}, \qquad a b + b a = 0, \qquad [ p, a ] = [ q, a ] = [ p, b ] = [ q, b ] = 0.
\]

\begin{enumerate}
\def\labelenumi{\alph{enumi})}
\setcounter{enumi}{2}
\tightlist
\item
  直到并包括问题 \(f\) ) 为止，我们都假设 \(\mathrm{E}_{p,q} = \mathrm{E}_{1 - p,1 - q} = \{0\}\) 。自同态 \(b\) 是单射；设 \(s\) 为从 \(\mathbf{R}\) 到 \(\mathbf{R}\) 的、满足下式的函数
\end{enumerate}

\[
s (x) = \left\{ \begin{array}{l l} 1 & \text { if } x > 0 \\ 0 & \text { if } x = 0 \\ - 1 & \text { if } x <   0. \end{array} \right.
\]

自同态 \(u = s(b)\) 是酉的。

\(d)\) 对每个 \(\varepsilon > 0\) ，我们有 \((|b| + \varepsilon)a = a(|b| + \varepsilon)\) 。

\begin{enumerate}
\def\labelenumi{\alph{enumi})}
\setcounter{enumi}{4}
\item
  我们有 \(uau^{-1} = -a\) （验证 \(s(b) = \lim_{\varepsilon \to 0} b(|b| + \varepsilon)^{-1}\) ）以及 \(ubu^{-1} = b\) 。
\item
  我们有 \(upu^{-1} = q\) 且 \(uqu^{-1} = p\) 。
\item
  我们假设 \(E_{p,q}\) 的希尔伯特维数等于 \(E_{1-p,1-q}\) 的希尔伯特维数。则存在 E 的酉自同态 u 使 \(upu^{-1} = q\) 且 \(uqu^{-1} = p\) 。（空间 \(E_{p,q}\) 与 \(E_{1-p,1-q}\) 正交；考虑希尔伯特空间 \(F = E_{p,q} \oplus E_{1-p,1-q}\) ，并对由 p 与 q 通过过渡到子空间导出的 \(F^{\circ}\) 的正交投影算子应用前面的情形。）
\item
  假设 \(\|p-q\|<1\) 。存在 E 的酉自同态 u 使 \(upu^{-1}=q\) 且 \(uqu^{-1}=p\) 。
\item
  假设 \(a = p - q\) 是紧的。自同态 \(k = qp\) 通过过渡到子空间诱导出弗雷德霍姆线性映射 \(\widetilde{k}\) （从 \(p\) 的像到 \(q\) 的像）。存在 E 的酉自同态 \(u\) 使 \(upu^{-1} = q\) 且 \(uqu^{-1} = p\) ，当且仅当 \(\widetilde{k}\) 的指标为零。
\end{enumerate}

\begin{enumerate}
\def\labelenumi{\arabic{enumi})}
\setcounter{enumi}{34}
\tightlist
\item
  设 E 为希尔伯特空间，u 为 E 上的正自伴部分算子。
\end{enumerate}

\begin{enumerate}
\def\labelenumi{\alph{enumi})}
\item
  对每个 \(t \in \mathbf{R}_{+}^{*}\) ，我们有 \(u \circ \mathrm{R}(u, -t) \in \mathcal{L}(\mathrm{E})\) 。
\item
  映射 \(f\) （从 \(\mathbf{R}_{+}^{*}\) 到 \(\mathcal{L}(\mathrm{E})\) ，由 \(f(t) = -t^{-1/2}u \circ \mathrm{R}(u, -t)\) 定义）关于 \(\mathbf{R}_{+}^{*}\) 上的勒贝格测度可积。
\item
  我们有
\end{enumerate}

\[
\sqrt {u} = \frac {1}{\pi} \int_ {\mathbf {R} _ {+} ^ {*}} f (t) d t.
\]

\begin{enumerate}
\def\labelenumi{\arabic{enumi})}
\setcounter{enumi}{35}
\item
  给出一个复希尔伯特空间上的正规部分算子的例子，其本质谱在 C 中不闭。
\item
  设 E 与 F 为复希尔伯特空间，u 为从 E 到 F 的稠定闭部分算子。若 u 是满射的，则存在 \(v \in \mathcal{L}(\mathrm{F}; \mathrm{E})\) 使 \(u \circ v = 1_{F}\) 。
\item
  设 u 为复希尔伯特空间 E 中的部分算子。
\end{enumerate}

\begin{enumerate}
\def\labelenumi{\alph{enumi})}
\item
  假设 u 是正的对称算子。设 \(F(u)\) 为 u 的弗里德里希斯扩张，q 为与 \(F(u)\) 相伴的正部分形式。证明 \(\mathrm{F}(u)\) 是 \(u\) 的唯一一个定义域包含于 \(\operatorname{dom}(q)\) 的自伴扩张。
\item
  假设 \(u\) 是具有稠密定义域的闭部分算子。设 \(q\) 为以 \(\operatorname{dom}(u)\) 为定义域、由 \(q(x,y) = \langle u(x) \mid u(y) \rangle\) 定义的正部分形式。证明表示 \(q\) 的自伴部分算子是 \(v = u^{*}u\) 。
\item
  我们假设 u 对称且 \(u^{2}\) 有稠密定义域。部分算子 \(u^{2}\) 对称且正；\(u^{2}\) 的弗里德里希斯扩张等于 \(u^{*}u\) 。
\end{enumerate}

\begin{enumerate}
\def\labelenumi{\arabic{enumi})}
\setcounter{enumi}{38}
\tightlist
\item
  设 \(n \geqslant 1\) 为整数。我们在 \(\mathbf{R}^n\) 上配备勒贝格测度。
\end{enumerate}

设 v 为 \(R^{n}\) 上局部平方可积的正函数。我们用 \(\Delta\) 表示 \(\mathcal{D}(\mathbf{R}^{n})\) 上或 \(\mathcal{D}'(\mathbf{R}^{n})\) 上的拉普拉斯算子，它由下式定义

\[
\Delta (f) = - \sum_ {i = 1} ^ {n} \partial_ {i} ^ {2} f.
\]

设 D 为 \(L^{2}(\mathbf{R}^{n})\) 上以 \(\mathcal{D}(\mathbf{R}^{n})\) 为定义域的部分算子，它满足

\[
\mathrm{D} (\varphi) = \Delta (\varphi + v \varphi).
\]

\begin{enumerate}
\def\labelenumi{\alph{enumi})}
\item
  部分算子 D 对称且正。
\item
  设 \(u \in \text{dom}(D^{*})\) 满足 \(D^{*}u + u = 0\) 。我们有 \(\Delta(u) + u = 0\) ，其中 \(u \in \text{L}^{2}(\mathbf{R}^{n})\) 的拉普拉斯算子是在广义函数意义下计算的。
\item
  广义函数 \(\Delta(u)\) 与 \((v+1)u\) 是局部可积函数，且广义函数 \(-\Delta(u)\) 是正的。（利用 IV 第 344 页习题 32。）
\item
  我们有 u = 0。（把 \(|u|\) 用无穷次可微函数 \(w_{\varepsilon}\) 通过卷积正则化，使得 \(\Delta(w_{\varepsilon}) \geqslant 0\) ，并应用上一问题。） e) 部分算子 D 是本质自伴的。
\end{enumerate}

\begin{enumerate}
\def\labelenumi{\arabic{enumi})}
\setcounter{enumi}{39}
\tightlist
\item
  设 E 为复希尔伯特空间。
\end{enumerate}

\begin{enumerate}
\def\labelenumi{\alph{enumi})}
\item
  设 u 为 E 上的正规部分算子，\((j, |u|)\) 为它的极分解。我们有 \(|u^{*}| = |u|\) ，且 \(u^{*}\) 的极分解是 \((j^{*}, |u|)\) 。我们有 \(j|u| = |u|j\) 。
\item
  设 u 为 E 上的自伴部分算子，\((j, |u|)\) 为它的极分解。我们有 \(j = \mathrm{s}(u)\) ，其中函数 s 由下式定义
\end{enumerate}

\[
\mathrm{s} (x) = \left\{ \begin{array}{l l} 1 & \text { if } x > 0 \\ 0 & \text { if } x = 0 \\ - 1 & \text { if } x <   0. \end{array} \right.
\]

\begin{enumerate}
\def\labelenumi{\alph{enumi})}
\setcounter{enumi}{2}
\tightlist
\item
  设 u 为 E 上的自伴部分算子，\((j, |u|)\) 为它的极分解。我们用 D 表示 \(|u|^{1/2}\) 的定义域，并通过令下式成立来定义以 D 为定义域的部分埃尔米特形式 q
\end{enumerate}

\[
q (x, y) = \langle | u | ^ {1 / 2} (x) \mid j | u | ^ {1 / 2} (y) \rangle .
\]

表示 \(q\) 的部分算子与 \(u\) 重合；若 \(u\) 是正的，则 \(q\) 是与 \(u\) 相伴的正部分形式。

\begin{enumerate}
\def\labelenumi{\alph{enumi})}
\setcounter{enumi}{3}
\tightlist
\item
  设 u 与 v 为 E 上的自伴部分算子。用 \(q_{u}\) 与 \(q_{v}\) 表示按问题 c) 的方式定义的部分埃尔米特形式。假设 \(\text{dom}(q_{u}) \cap \text{dom}(q_{v})\) 在 E 中稠密。u 与 v 在埃尔米特形式意义下的和是 E 上表示部分形式的部分算子 \(u + v\)
\end{enumerate}

\[
q (x, y) = q _ {u} (x, y) + q _ {v} (x, y)
\]

其定义域为 \(\mathrm{dom}(q_{u})\cap\mathrm{dom}(q_{v})\) 。证明部分算子 \(u+v\) 是 \(u+v\) 的扩张；若 u 与 v 是正的，则 \(u+v\) 是正的自伴部分算子。

\begin{enumerate}
\def\labelenumi{\alph{enumi})}
\setcounter{enumi}{4}
\item
  给出一个例子，其中 \(u + v\) 的定义域不稠密，但 \(u + v\) 有定义。
\item
  设 \(n \geqslant 1\) 为整数，\(S \subset R^{n}\) 为勒贝格测度为零的闭集。设 v 为 \(R^{n}\) 上关于勒贝格测度可测的正函数，且 v 在 \(R^{n}-S\) 上局部可积。自伴部分算子 \(\Delta + m_{v}\) 有定义。
\end{enumerate}

\begin{enumerate}
\def\labelenumi{\arabic{enumi})}
\setcounter{enumi}{40}
\tightlist
\item
  设 \(n \geqslant 1\) 为整数。我们用 \(\Delta\) 表示 \(\mathbf{R}^n\) 上的拉普拉斯算子，用 \(u\) 表示 \(\Delta\) 到子空间 \(\mathcal{D}(\mathbf{R}^n - \{0\})\) 的约化。
\end{enumerate}

\begin{enumerate}
\def\labelenumi{\alph{enumi})}
\item
  u 的定义域在希尔伯特空间 \(E_{q}\) （它同与 \(\Delta\) 相伴的部分形式 q 相伴）中稠密，当且仅当 \(n \geqslant 2\) 。（把 \(E_{q}\) 的对偶与一个缓增广义函数空间等同起来；若 \(g \in E_{q}'\) 在 dom(u) 上为零，证明 g 是支撑包含于 \(\{0\}\) 的广义函数；然后利用 IV 第 338 页习题 24 得出结论。）
\item
  \(u\) 的定义域在希尔伯特空间 \(\mathrm{E}_{\Delta}\) 中稠密，当且仅当 \(n \geqslant 4\) 。
\end{enumerate}

\begin{enumerate}
\def\labelenumi{\arabic{enumi})}
\setcounter{enumi}{41}
\tightlist
\item
  设 E 为复希尔伯特空间，u 为 E 上具有稠密定义域的闭部分算子。我们称 E 上的部分算子 v 关于 u 是相对紧的，如果 v 的定义域包含 u 的定义域，且 v 到 \(\text{dom}(u)\) 的限制是 \(E_{u}\) 到 E 的紧线性映射。
\end{enumerate}

\begin{enumerate}
\def\labelenumi{\alph{enumi})}
\item
  每个部分算子 \(v\) （关于 \(u\) 相对紧）都关于 \(u\) 有界，且其相对范数 \(\| v \|_u\) 为零。（若相对范数不为零，构造收敛于 0 的序列 \((x_n)\) （在 \(\text{dom}(u)\) 中），使 \(\| v(x_n) \| = 1\) （对所有 \(n\) ），且使 \((u(x_n))\) 弱收敛于 E 中的元素 \(y\) ；证明 \(y = 0\) ，由此推出 \(v\) 不是相对紧的。）
\item
  设 \(\lambda \in \mathbf{C} - \mathrm{Sp}(u)\) 。设 \(v\) 为 E 上的部分算子，其定义域包含 \(\mathrm{dom}(u)\) 。则 \(v\) 关于 \(u\) 相对紧，当且仅当 \(v \circ \mathrm{R}(u, \lambda)\) 是 E 的紧自同态。
\end{enumerate}

\begin{enumerate}
\def\labelenumi{\arabic{enumi})}
\setcounter{enumi}{42}
\tightlist
\item
  设 E 为复希尔伯特空间。
\end{enumerate}

\begin{enumerate}
\def\labelenumi{\alph{enumi})}
\tightlist
\item
  设 u 为 E 上的自伴部分算子。对每个 \(\lambda \in R\) ，我们有 \(\lambda \in \mathrm{Sp}_{\mathrm{e}}(u)\) ，当且仅当存在 E 中的序列 \((x_{n})\) ，它在 E 中弱收敛于 0 且满足条件
\end{enumerate}

\[
\liminf _ {n \to + \infty} \| x _ {n} \| > 0, \quad \lim _ {n \to + \infty} (u (x _ {n}) - \lambda x _ {n}) = 0.
\]

\begin{enumerate}
\def\labelenumi{\alph{enumi})}
\setcounter{enumi}{1}
\item
  设 \(u\) 与 \(v\) 为 E 上的自伴部分算子。若存在 \(\lambda \in \mathbf{C} - (\mathrm{Sp}(u) \cup \mathrm{Sp}(v))\) 使 \(\mathrm{R}(u, \lambda) - \mathrm{R}(v, \lambda)\) 紧，则 \(\mathrm{Sp}_{\mathrm{e}}(u) = \mathrm{Sp}_{\mathrm{e}}(v)\) 。
\item
  设 u 为 E 上的闭对称部分算子。设 \(u_{1}\) 与 \(u_{2}\) 为 u 的自伴扩张。若空间 \(\operatorname{Ker}(u^{*}-i)\) 是有限维的，则 \(\operatorname{Sp}_{\mathrm{e}}(u_{1})=\operatorname{Sp}_{\mathrm{e}}(u_{2})\) 。
\item
  设 u 为 E 上的自伴部分算子，v 为关于 u 相对紧的部分算子（见上一习题）。部分算子 \(u + v\) 自伴，且 \(\mathrm{Sp}_{\mathrm{e}}(u + v) = \mathrm{Sp}_{\mathrm{e}}(u)\) 。
\end{enumerate}

